% Fragmento público generado desde propietarios íntegros.
% Residencia: 52.10.1.
% Fuente: ../incoming/radion_monografia_20260827/manuscrito/sections/11_modularidad_orbitas.tex:1-84
\noindent Los caracteres modulares, orbitales y de camino conservan invariantes distintos de una misma forma orientada.\par\medskip

\subsubsection{Toro complejo asociado a la elipse de acción}

\paragraph{Retículo rectangular de períodos}

La elipse por sí sola no es un toro. Declaramos el retículo
\begin{equation}
\Lambda_\eta=a_S\Z+i,b_S\Z
\label{eq:lattice}
\end{equation}
y el cociente \(\C/\Lambda_\eta\). Su módulo es
\begin{equation}
\tau_\eta=i\frac{b_S}{a_S}=ie^{-\eta}
=i\,0.741048144159375160795483663369300\ldots.
\label{eq:tau-eta}
\end{equation}
La inversión bidireccional \(\eta\mapsto-\eta\) intercambia los ejes:
\begin{equation}
\tau_{-\eta}=ie^\eta=-\frac1{\tau_\eta}.
\label{eq:modular-S}
\end{equation}
Es exactamente la transformación modular \(S:\tau\mapsto-1/\tau\).

\paragraph{Invariancia de la función modular \(j\)}

Por invariancia modular,
\begin{equation}
\boxed{j(\tau_{-\eta})=j(\tau_\eta).}
\label{eq:j-invariance}
\end{equation}
En la normalización \(j(i)=1728\),
\[
j(\tau_\eta)
=5597.43843932408729830097089254768\ldots.
\]
La función moonshine normalizada \(J=j-744\) da
\[
J(\tau_\eta)
=4853.43843932408729830097089254768\ldots.
\]
Para el bloque APP basal,
\[
\tau_0=i\frac{\sqrt5}{3},
\qquad
j(\tau_0)=5369.48832024674306021916882626352\ldots.
\]

\begin{proposition}[Lo que conserva la clase modular]
La involución de hoja invierte \(\eta\), pero \(j\) conserva la clase del
toro. Los focos y el etiquetado de ejes retienen la orientación que \(j\)
olvida.
\end{proposition}

Este resultado conecta de manera exacta la doble vía con la modularidad una
vez declarado el cociente. No autoriza a hacer actuar directamente al
Monstruo sobre la carga, los dígitos o los focos. El corredor excepcional
\[
\mathrm{APP}\to\mathrm{Paley}\to\mathrm{Witt}\to\mathrm{Golay}
\to A_2^{12}\to\mathrm{Leech}\to V^\natural
\to\mathrm{Monster}\to\mathrm{Moonshine}
\]
es otra proyección del mismo estado enriquecido. Para identificar una acción
concreta sobre el radión se necesita un entrelazador que preserve sus
coordenadas pertinentes.

\paragraph{Cadena de invariantes del carácter de forma}

Reuniendo las realizaciones declaradas,
\begin{equation}
\boxed{
r_\eta=e^{-\eta}
=\sqrt{\frac{A-C^*}{A+C^*}}
=\frac{b_1}{a_1}
=\frac{b_S}{a_S}
=\frac{\Delta P}{\Delta Q}
=\frac{Z_\eta}{Z_*}
=\Im\tau_\eta.}
\label{eq:shape-master}
\end{equation}
Una sola razón recorre geometría espectral, acción, covarianza, impedancia y
módulo. Cada igualdad se apoya en un mapa explícito; ninguna depende de una
semejanza decimal.

